% The five senses in which the literature uses the word "skill".
\begin{table*}[t]
\centering
\footnotesize
\renewcommand{\arraystretch}{1.25}
\begin{tabular}{@{}p{0.16\textwidth} p{0.22\textwidth} p{0.18\textwidth} c c c c@{}}
\toprule
\textbf{Sense of ``skill''} & \textbf{Representation} & \textbf{Examples} &
\textbf{Inspect.} & \textbf{Adapt.} & \textbf{Compos.} & \textbf{Distrib.} \\
\midrule
latent policy & latent-conditioned network $\pi(a\mid s,z)$ & DIAYN, DADS, METRA & \nn & \pp & \nn & \nn \\
option / primitive & temporally-extended sub-policy with a skill prior & SPiRL, OPAL, SkiMo & \pp & \pp & \yy & \nn \\
code & executable program over an API & Code-as-Policies, Voyager & \yy & \yy & \yy & \pp \\
robot app & packaged, one-tap deployable behaviour & UniStore motion packages & \nn & \nn & \nn & \yy \\
market product & a listed, versioned, priced commodity & UniStore listings & \nn & \nn & \pp & \yy \\
\bottomrule
\end{tabular}
\caption{The word ``skill'' spans at least five senses that the literature frequently
conflates. Columns record whether a skill is human-\textbf{inspect}able, \textbf{adapt}able
after creation, symbolically \textbf{compos}able, and \textbf{distrib}utable across robots
(\yy\ yes, \pp\ partial, \nn\ no). Only the \emph{code} sense is inspectable, adaptable, and
composable at once; only the app and market senses are readily distributable, which is the
mismatch the skill economy (\S\ref{sec:open}) must close.}
\label{tab:skill_senses}
\end{table*}
